\documentclass[11pt]{article}
\usepackage[margin=1in]{geometry}
\usepackage{amsmath,amssymb,amsthm}
\usepackage{graphicx}
\usepackage[hidelinks]{hyperref}

\newtheorem{theorem}{Theorem}
\newtheorem{lemma}{Lemma}
\newtheorem{corollary}{Corollary}
\newtheorem{remark}{Remark}

\newcommand{\C}{\mathbb{C}}
\newcommand{\N}{\mathbb{N}}
\newcommand{\R}{\mathbb{R}}
\newcommand{\Id}{\mathrm{Id}}
\newcommand{\norm}[1]{\left\lVert #1 \right\rVert}
\newcommand{\abs}[1]{\left\lvert #1 \right\rvert}
\newcommand{\Sigmaa}[1]{\Sigma_{#1}}
\newcommand{\eps}{\varepsilon}
\newcommand{\Res}{\mathrm{R}}

\title{Removing the Extra Smoothing Factor in Fackler's Kato--Beurling Converse}
\author{Agent lane 14}
\date{July 18, 2026}

\begin{document}
\maketitle

\section*{Source and Target}

The source paper is Stephan Fackler, \emph{Regularity of Semigroups via the
Asymptotic Behaviour at Zero}, arXiv:1203.5498.  Corollary 2.5 proves that if
a \(C_0\)-semigroup \((T(t))_{t\geq0}\) is holomorphic, then for all large
\(N\) one has
\[
  \limsup_{t\downarrow0}
  \norm{(T(t)-\Id)^N T(t)^K}^{1/N}<2
\]
after choosing a sufficiently large integer \(K\).  Remark 2.6 observes that,
for a bounded holomorphic semigroup on a Hilbert space whose negative
generator has a bounded \(H^\infty\)-calculus, the extra factor \(T(t)^K\)
can be omitted, and asks whether this remains true in general.

\begin{center}
\includegraphics[width=0.95\textwidth]{figures/open_problem_crop.png}
\end{center}

This packet gives an affirmative answer.  The proof does not use Hilbert
space structure or an \(H^\infty\)-calculus.

\section*{Proof Intuition}

For a bounded holomorphic semigroup \(T(t)=e^{-tA}\), the negative generator
\(A\) is sectorial of some angle strictly smaller than \(\pi/2\).  Choose a
larger sector angle \(\psi<\pi/2\).  On the two boundary rays of this sector,
the scalar function \(z\mapsto e^{-z}-1\) is uniformly bounded by a number
\(q_\psi<2\).  Thus the scalar finite-difference function
\[
  f_N(z)=(e^{-z}-1)^N
\]
has boundary size at most \(q_\psi^N\), apart from its harmless limit
\((-1)^N\) at infinity.  A sectorial contour estimate then gives
\[
  \norm{f_N(tA)} \le 1+C\log(N+1)q_\psi^N+o(1)
\]
uniformly for \(t>0\).  Since \(q_\psi<2\), the \(N\)-th root is eventually
strictly less than \(2\).  Finally, an arbitrary holomorphic \(C_0\)-semigroup
can be exponentially rescaled to a bounded holomorphic semigroup on a smaller
sector, and this rescaling does not change the small-time \(N\)-th finite
difference limsup.

\section*{Sectorial Estimate}

We use the standard notation
\[
  \Sigma_\alpha=\{z\in\C\setminus\{0\}: \abs{\arg z}<\alpha\}.
\]
Recall that if \(-A\) generates a bounded holomorphic \(C_0\)-semigroup, then
\(A\) is sectorial of some angle \(\theta<\pi/2\).

\begin{lemma}[A boundary estimate for finite-limit functions]
\label{lem:sectorial-estimate}
Let \(A\) be sectorial of angle \(\theta<\pi/2\), and fix
\(\psi\in(\theta,\pi/2)\).  There is a constant \(C_{\psi,A}\) with the
following property.  Let \(f\) be holomorphic in a neighbourhood of
\(\overline{\Sigma_\psi}\), continuous at \(0\) and at infinity, with
\(f(0)=0\) and \(f(\infty)=L\).  Put
\[
  m_f(r)=\max_{\sigma=\pm1}\abs{f(re^{i\sigma\psi})},\qquad
  m_f^\infty(r)=\max_{\sigma=\pm1}\abs{f(re^{i\sigma\psi})-L}.
\]
Then for every \(R\geq 1\),
\[
  \norm{f(A)}
  \le
  \abs{L}
  + C_{\psi,A}\left(
      \int_0^1 m_f(r)\,\frac{dr}{r}
      +\int_1^R m_f(r)\,\frac{dr}{r}
      +\int_R^\infty m_f^\infty(r)\,\frac{dr}{r}
    \right).
\]
\end{lemma}

\begin{proof}
This is the elementary sectorial functional-calculus estimate obtained from
the usual keyhole contour around \(\Sigma_\theta\).  For completeness we spell
out the point of the estimate.  Let
\[
  M_{\psi,A}:=\sup_{\lambda\notin\Sigma_\psi}
       \norm{\lambda(\lambda-A)^{-1}}<\infty .
\]
On the two rays \(\lambda=re^{\pm i\psi}\), the resolvent contributes at most
\(M_{\psi,A}\,dr/r\).  Since \(f(0)=0\), the small circular arc vanishes as its
radius tends to \(0\).  On the large arc, the constant part \(L\) contributes
exactly \(L\Id\), and the remainder is controlled by \(m_f^\infty\).  Taking
the inner radius to \(0\) and the outer radius to infinity gives the displayed
bound.  This is the standard extended sectorial calculus for functions with
finite limits at \(0\) and \(\infty\), with the norm estimate read directly
from the same contour.
\end{proof}

\begin{lemma}[Uniform finite-difference bound]
\label{lem:finite-difference-bound}
Let \(-A\) generate a bounded holomorphic \(C_0\)-semigroup on a Banach space.
Then there exist \(q<2\) and \(C>0\) such that for every \(N\geq2\) and every
\(t>0\),
\[
  \norm{(e^{-tA}-\Id)^N}
  \le
  1+C\left(\frac1N+\log(N+1)q^N+e^{-N}\right).
\]
Consequently,
\[
  \limsup_{N\to\infty}\ \sup_{t>0}
  \norm{(e^{-tA}-\Id)^N}^{1/N}<2.
\]
\end{lemma}

\begin{proof}
Choose \(\psi\in(\theta,\pi/2)\), where \(\theta\) is a sectoriality angle for
\(A\), and put
\[
  q:=\sup_{r>0}\abs{1-e^{-re^{i\psi}}}.
\]
Then \(q<2\): the displayed function is continuous in \(r\), tends to \(0\) at
\(r=0\), tends to \(1\) as \(r\to\infty\), and cannot equal \(2\) for finite
\(r\) because \(\abs{e^{-re^{i\psi}}}=e^{-r\cos\psi}<1\).

Fix \(N\ge2\) and set \(f_N(z)=(e^{-z}-1)^N\).  The operator
\((e^{-tA}-\Id)^N\) is \(f_N(tA)\).  The scaled operator \(tA\) has the same
sectorial angle and the same resolvent bound \(M_{\psi,A}\), so Lemma
\ref{lem:sectorial-estimate} applies with constants independent of \(t\).
Since \(f_N(0)=0\) and \(f_N(\infty)=(-1)^N\), it remains only to estimate the
three scalar integrals in Lemma \ref{lem:sectorial-estimate}.

For \(0<r\le1\), the inequality
\[
  \abs{1-e^{-w}}\le \abs{w}\qquad(\mathrm{Re}\,w\ge0)
\]
gives \(m_{f_N}(r)\le r^N\), and hence
\[
  \int_0^1 m_{f_N}(r)\,\frac{dr}{r}\le \frac1N.
\]
For \(1\le r\le R\), \(m_{f_N}(r)\le q^N\), so
\[
  \int_1^R m_{f_N}(r)\,\frac{dr}{r}\le q^N\log R.
\]
Finally let \(c=\cos\psi>0\) and choose \(R=N/c\).  If \(r\ge R\) and
\(u=e^{-re^{i\psi}}\), then \(\abs{u}\le e^{-N}\), and
\[
  \abs{(1-u)^N-1}
  \le N\abs{u}(1+\abs{u})^{N-1}
  \le 2N e^{-cr}.
\]
The same estimate holds on the lower boundary ray.  Thus
\[
  \int_R^\infty m_{f_N}^\infty(r)\,\frac{dr}{r}
  \le 2N\int_R^\infty e^{-cr}\,\frac{dr}{r}
  \le 2e^{-N}.
\]
Substitution into Lemma \ref{lem:sectorial-estimate} gives the asserted
uniform bound after increasing \(C\).  Since \(q<2\), the \(N\)-th root of the
right hand side is eventually strictly smaller than \(2\).
\end{proof}

\section*{Main Result}

\begin{theorem}
\label{thm:main}
Let \((T(t))_{t\ge0}\) be a holomorphic \(C_0\)-semigroup on a Banach space
\(X\).  Then there exists \(N\in\N\) such that
\[
  \limsup_{t\downarrow0}
  \norm{(T(t)-\Id)^N}^{1/N}<2.
\]
In particular, the factor \(T(t)^K\) in Fackler's Corollary 2.5(b) can be
omitted.
\end{theorem}

\begin{proof}
First suppose that \((T(t))\) is bounded holomorphic, and write
\(T(t)=e^{-tA}\).  Lemma \ref{lem:finite-difference-bound} gives a sequence
of integers \(N\) for which
\[
  \sup_{t>0}\norm{(T(t)-\Id)^N}^{1/N}<2,
\]
and hence the required small-time limsup estimate.

For a general holomorphic \(C_0\)-semigroup, restrict the holomorphic
extension to a slightly smaller sector.  By the semigroup law and local
boundedness near \(0\), there is \(\omega\in\R\) such that
\[
  S(z):=e^{-\omega z}T(z)
\]
is bounded holomorphic on that smaller sector.  The bounded case gives an
integer \(N\) with
\[
  \limsup_{t\downarrow0}\norm{(S(t)-\Id)^N}^{1/N}<2.
\]
Since \(T(t)=e^{\omega t}S(t)\), we have
\[
  T(t)-\Id=e^{\omega t}(S(t)-\Id)+(e^{\omega t}-1)\Id .
\]
For this fixed \(N\), expanding the \(N\)-th power shows that
\[
  (T(t)-\Id)^N-e^{N\omega t}(S(t)-\Id)^N\to0
\]
in operator norm as \(t\downarrow0\), because \((S(t))\) is bounded near
\(0\).  As \(e^{N\omega t}\to1\), the desired strict inequality transfers
from \(S\) to \(T\).
\end{proof}

\section*{Verification and Novelty Notes}

No numerical computation is used.  The proof is an analytic sectorial-calculus
argument.  The main review point is Lemma \ref{lem:sectorial-estimate}; it is
the standard contour estimate for sectorial operators, but it is where all
operator-theoretic content is concentrated.

Local duplicate checks were run against
\texttt{runs/fa\_banach\_001/registry\_index.tsv},
\texttt{solutions/index.tsv}, \texttt{attempts/index.tsv}, and
\texttt{proof\_gaps/index.tsv} for arXiv id \texttt{1203.5498}, the phrases
\texttt{additional factor}, \texttt{z\^{}K}, \texttt{Kato-Beurling}, and the
finite-difference expressions in Corollary 2.5.  No completed existing packet
for this exact question was found before this promotion.

A bounded web search on July 18, 2026 used the exact phrases
\texttt{additional factor z\^{}K Kato-Beurling},
\texttt{limsup (T(t)-Id)\^{}N T(t)\^{}K < 2},
\texttt{Pazy limsup T(t)-I holomorphic counterexample}, and the source paper
title.  It found adjacent literature, including Batty--Gomilko--Tomilov's
Besov-calculus paper and Hytönen--Lappas's quantitative Kato estimates for
bounded holomorphic semigroups, but no separate paper explicitly answering
Fackler's Remark 2.6 by removing the \(z^K\) factor.  Novelty confidence is
\texttt{bounded-search likely new}; this is not a full bibliographic audit.

\section*{Limitations and Review Recommendation}

The theorem answers the question as stated for holomorphic \(C_0\)-semigroups
on arbitrary Banach spaces.  It does not provide an explicit optimal \(N\);
the proof gives existence through the sectoriality angle and resolvent
constant.

Human review should check:
\begin{itemize}
  \item the contour estimate in Lemma \ref{lem:sectorial-estimate}, especially
  the treatment of the finite limit \(f(\infty)\);
  \item the transfer from a general holomorphic semigroup to an exponentially
  rescaled bounded holomorphic semigroup;
  \item the scalar tail estimate for
  \((1-e^{-re^{i\psi}})^N-1\).
\end{itemize}

\begin{thebibliography}{9}
\bibitem{Fackler}
S. Fackler,
\emph{Regularity of Semigroups via the Asymptotic Behaviour at Zero},
Semigroup Forum 87 (2013), 21--42; arXiv:1203.5498.

\bibitem{Haase}
M. Haase,
\emph{The Functional Calculus for Sectorial Operators},
Operator Theory: Advances and Applications 169, Birkhauser, 2006.

\bibitem{BattyGomilkoTomilov}
C. J. K. Batty, A. Gomilko, and Y. Tomilov,
\emph{A Besov algebra calculus for generators of operator semigroups and
related norm-estimates},
Mathematische Annalen 379 (2021), 23--93.

\bibitem{HytonenLappas}
T. Hytönen and S. Lappas,
\emph{Quantitative estimates for bounded holomorphic semigroups},
Semigroup Forum 108 (2024), 115--144.
\end{thebibliography}

\end{document}
